\bookchapter{Lab: Rebuild Everything from Code}{ch:11-lab-rebuild-everything-from-code}

\begin{chapterintro}

In this lab Terraform builds a complete small system: tinyapp, a PostgreSQL database, a private network between them, and a volume for the data. Then you change it, ``lose the server'', get it all back with one command, and finally destroy and recreate everything, learning on the way what code can rebuild and what it can't.

\end{chapterintro}

\emph{The problem.} My server died. Can I get it all back in five minutes?

\emph{The question.} If everything is described in code, can I delete it all and rebuild it with one command?

\section{What You'll Build}\label{11-lab-rebuild-everything-from-code:what-youll-build}

\codeneed{12}\begin{codeblock}
\begin{Verbatim}[breaklines=false, fontsize=\fontsize{8.8}{8.68}\selectfont, fontfamily=diagrammono, formatcom=\diagramlines]
            your machine, port 8090
                     │
 ┌───────────────────┼─ network: tinyapp-net ──────────────┐
 │                   ▼                                     │
 │  ┌──────────────────────┐      ┌─────────────────────┐  │
 │  │ tinyapp-web          │ ───► │ tinyapp-db          │  │
 │  │ localhost:5001/      │ 5432 │ postgres:16-alpine  │  │
 │  │   tinyapp:1.0.0      │      └──────────┬──────────┘  │
 │  └──────────────────────┘                 │             │
 └───────────────────────────────────────────┼─────────────┘
                                             ▼
                                volume: tinyapp-db-data
\end{Verbatim}
\end{codeblock}

Six resources: two images, two containers, one network, one volume. The app image is the \texttt{1.0.0} you pushed in Chapter 9, so the registry must be running. If \texttt{curl\ localhost:}\allowbreak{}\texttt{5001/}\allowbreak{}\texttt{v2/}\allowbreak{}\texttt{tinyapp/}\allowbreak{}\texttt{tags/}\allowbreak{}\texttt{list} doesn't list \texttt{1.0.0}, build and push it from the book's \texttt{labs/tinyapp/} folder with \texttt{docker\ build\ -}\allowbreak{}\texttt{t\ localhost:}\allowbreak{}\texttt{5001/}\allowbreak{}\texttt{tinyapp:}\allowbreak{}\texttt{1.}\allowbreak{}\texttt{0.}\allowbreak{}\texttt{0\ .} and \texttt{docker\ push\ localhost:}\allowbreak{}\texttt{5001/}\allowbreak{}\texttt{tinyapp:}\allowbreak{}\texttt{1.}\allowbreak{}\texttt{0.}\allowbreak{}\texttt{0}. This system uses port 8090 and its own names, so it doesn't disturb Chapter 10's \texttt{tinyapp} on port 8080.

\codeneed{10}

\section{Step 1: Describe the System}\label{11-lab-rebuild-everything-from-code:step-1-describe-the-system}

Create \texttt{\textasciitilde{}/}\allowbreak{}\texttt{infra-}\allowbreak{}\texttt{labs/}\allowbreak{}\texttt{ch11/}\allowbreak{}\texttt{main.}\allowbreak{}\texttt{tf}. It starts like Chapter 8's file, then declares two inputs:

\codeneed{4}\begin{codeblock}

\begin{Shaded}
\begin{Highlighting}[]
\NormalTok{terraform \{}
\NormalTok{  required\_providers \{}
\NormalTok{    docker = \{}
\NormalTok{      source  = "kreuzwerker/docker"}
\NormalTok{      version = "\textasciitilde{}\textgreater{} 4.6"}
\NormalTok{    \}}
\NormalTok{  \}}
\NormalTok{\}}

\NormalTok{provider "docker" \{\}}

\NormalTok{variable "app\_version" \{}
\NormalTok{  type    = string}
\NormalTok{  default = "1.0.0"}
\NormalTok{\}}

\NormalTok{variable "db\_password" \{}
\NormalTok{  type      = string}
\NormalTok{  sensitive = true}
\NormalTok{\}}
\end{Highlighting}
\end{Shaded}

\end{codeblock}

A \textbf{variable}\index{variable} is an input to a configuration. \texttt{app\_version} has a default. \texttt{db\_password} has none, so Terraform must be given it, and \texttt{sensitive\ =\ true} stops Terraform printing it. Terraform reads any environment variable called \texttt{TF\_VAR\_\textless{}name\textgreater{}}, so the password never has to be written in a file.

\codeneed{7}

Next, the network, the volume, and the two images:

\codeneed{4}\begin{codeblock}

\begin{Shaded}
\begin{Highlighting}[]
\NormalTok{resource "docker\_network" "lab" \{}
\NormalTok{  name = "tinyapp{-}net"}
\NormalTok{\}}

\NormalTok{resource "docker\_volume" "db\_data" \{}
\NormalTok{  name = "tinyapp{-}db{-}data"}
\NormalTok{\}}

\NormalTok{resource "docker\_image" "postgres" \{}
\NormalTok{  name         = "postgres:16{-}alpine"}
\NormalTok{  keep\_locally = true}
\NormalTok{\}}

\NormalTok{resource "docker\_image" "app" \{}
\NormalTok{  name         = "localhost:5001/tinyapp:1.0.0"}
\NormalTok{  keep\_locally = true}
\NormalTok{\}}
\end{Highlighting}
\end{Shaded}

\end{codeblock}

\codeneed{7}

The database container joins the network, keeps its files in the volume, and has a health check. \texttt{wait\ =\ true} makes Terraform wait until the container is healthy before moving on:

\codeneed{4}\begin{codeblock}

\begin{Shaded}
\begin{Highlighting}[]
\NormalTok{resource "docker\_container" "db" \{}
\NormalTok{  name  = "tinyapp{-}db"}
\NormalTok{  image = docker\_image.postgres.image\_id}
\NormalTok{  env = [}
\NormalTok{    "POSTGRES\_USER=tinyapp",}
\NormalTok{    "POSTGRES\_PASSWORD=$\{var.db\_password\}",}
\NormalTok{    "POSTGRES\_DB=tinyapp",}
\NormalTok{  ]}

\NormalTok{  networks\_advanced \{}
\NormalTok{    name = docker\_network.lab.name}
\NormalTok{  \}}

\NormalTok{  volumes \{}
\NormalTok{    volume\_name    = docker\_volume.db\_data.name}
\NormalTok{    container\_path = "/var/lib/postgresql/data"}
\NormalTok{  \}}

\NormalTok{  healthcheck \{}
\NormalTok{    test     = ["CMD", "pg\_isready", "{-}U", "tinyapp"]}
\NormalTok{    interval = "2s"}
\NormalTok{    retries  = 10}
\NormalTok{  \}}
\NormalTok{  wait = true}
\NormalTok{\}}
\end{Highlighting}
\end{Shaded}

\end{codeblock}

\codeneed{7}

Finally, the app and an output:

\codeneed{4}\begin{codeblock}

\begin{Shaded}
\begin{Highlighting}[]
\NormalTok{resource "docker\_container" "app" \{}
\NormalTok{  name  = "tinyapp{-}web"}
\NormalTok{  image = docker\_image.app.image\_id}
\NormalTok{  env = [}
\NormalTok{    "APP\_VERSION=$\{var.app\_version\}",}
\NormalTok{    "DATABASE\_URL=postgresql://tinyapp:$\{var.db\_password\}@$\{docker\_container.db.name\}:5432/tinyapp",}
\NormalTok{  ]}

\NormalTok{  networks\_advanced \{}
\NormalTok{    name = docker\_network.lab.name}
\NormalTok{  \}}

\NormalTok{  ports \{}
\NormalTok{    internal = 8000}
\NormalTok{    external = 8090}
\NormalTok{  \}}
\NormalTok{  wait = true}
\NormalTok{\}}

\NormalTok{output "url" \{}
\NormalTok{  value = "http://localhost:8090"}
\NormalTok{\}}
\end{Highlighting}
\end{Shaded}

\end{codeblock}

\texttt{\$\{docker\_}\allowbreak{}\texttt{container.}\allowbreak{}\texttt{db.}\allowbreak{}\texttt{name\}} is a reference inside the database address. It gives the app the database's name on the network, and it tells Terraform to create the database first. \texttt{wait\ =\ true} here waits for the image's own \texttt{HEALTHCHECK} from Chapter 9. An \textbf{output}\index{output} is a value Terraform prints after \texttt{apply} and shows again with \texttt{terraform\ output}.

\section{Step 2: Build It}\label{11-lab-rebuild-everything-from-code:step-2-build-it}

\codeneed{17}\begin{codeblock}

\begin{Shaded}
\begin{Highlighting}[]
\BuiltInTok{cd}\NormalTok{ \textasciitilde{}/infra{-}labs/ch11}
\BuiltInTok{export} \VariableTok{TF\_VAR\_db\_password}\OperatorTok{=}\NormalTok{change{-}me{-}locally}
\ExtensionTok{terraform}\NormalTok{ init}
\ExtensionTok{terraform}\NormalTok{ plan}
\end{Highlighting}
\end{Shaded}

\end{codeblock}

\codeneed{12}\begin{codeblock}
\begin{Verbatim}[formatcom=\outputstyle]
…
  # docker_container.app will be created
  # docker_container.db will be created
  # docker_image.app will be created
  # docker_image.postgres will be created
  # docker_network.lab will be created
  # docker_volume.db_data will be created
…
Plan: 6 to add, 0 to change, 0 to destroy.

Changes to Outputs:
  + url = "http://localhost:8090"
\end{Verbatim}
\end{codeblock}

\codeneed{17}

(These lines are picked out of a longer plan.) Apply it, answering \texttt{yes}. Putting \texttt{time} in front measures how long it takes:

\codeneed{14}\begin{codeblock}

\begin{Shaded}
\begin{Highlighting}[]
\BuiltInTok{time}\NormalTok{ terraform apply}
\end{Highlighting}
\end{Shaded}

\end{codeblock}

\codeneed{12}\begin{codeblock}
\begin{Verbatim}[formatcom=\outputstyle]
…
docker_network.lab: Creation complete after 2s [id=309f6b4865e3…]
docker_container.db: Creating...
docker_container.db: Creation complete after 4s [id=0506b66e5fb6…]
docker_container.app: Creating...
docker_container.app: Creation complete after 6s [id=19b358f8dfa7…]

Apply complete! Resources: 6 added, 0 changed, 0 destroyed.

Outputs:

url = "http://localhost:8090"
\end{Verbatim}
\end{codeblock}

It took about 12 seconds; \texttt{time} adds its own lines, which differ between shells. Notice the order: Terraform created the network, volume, and images together, then the database, and only when the database was healthy, the app. Try it:

\codeneed{7}\begin{codeblock}

\begin{Shaded}
\begin{Highlighting}[]
\ExtensionTok{curl}\NormalTok{ localhost:8090/health}
\ExtensionTok{curl}\NormalTok{ localhost:8090/visits}
\ExtensionTok{curl}\NormalTok{ localhost:8090/visits}
\end{Highlighting}
\end{Shaded}

\end{codeblock}

\codeneed{3}\begin{codeblock}
\begin{Verbatim}[formatcom=\outputstyle]
{"status":"ok","version":"1.0.0"}
{"visits":1}
{"visits":2}
\end{Verbatim}
\end{codeblock}

The count is now stored in PostgreSQL.

\codeneed{11}

\section{Step 3: Change Something}\label{11-lab-rebuild-everything-from-code:step-3-change-something}

In \texttt{main.tf}, change the default of \texttt{app\_version} from \texttt{"1.0.0"} to \texttt{"1.1.0"}, and plan:

\codeneed{5}\begin{codeblock}
\begin{Verbatim}[formatcom=\outputstyle]
  # docker_container.app must be replaced
…
      ~ env                                         = (sensitive value) # forces replacement
…
Plan: 1 to add, 0 to change, 1 to destroy.
\end{Verbatim}
\end{codeblock}

Only the app container changes. Its environment is shown as \texttt{(sensitive\ value)} because it contains the password. Apply it, and Terraform destroys and re-creates just that container.

\codeneed{5}\begin{codeblock}

\begin{Shaded}
\begin{Highlighting}[]
\ExtensionTok{curl}\NormalTok{ localhost:8090/health}
\ExtensionTok{curl}\NormalTok{ localhost:8090/visits}
\end{Highlighting}
\end{Shaded}

\end{codeblock}

\codeneed{2}\begin{codeblock}
\begin{Verbatim}[formatcom=\outputstyle]
{"status":"ok","version":"1.1.0"}
{"visits":3}
\end{Verbatim}
\end{codeblock}

A new app container, and the count carried on, because it lives in the database, not in the app.

\codeneed{8}

\section{Step 4: The Server Dies}\label{11-lab-rebuild-everything-from-code:step-4-the-server-dies}

Simulate losing the machine's containers and network all at once:

\codeneed{2}\begin{codeblock}

\begin{Shaded}
\begin{Highlighting}[]
\ExtensionTok{docker}\NormalTok{ rm }\AttributeTok{{-}f}\NormalTok{ tinyapp{-}web tinyapp{-}db}
\ExtensionTok{docker}\NormalTok{ network rm tinyapp{-}net}
\end{Highlighting}
\end{Shaded}

\end{codeblock}

\texttt{curl\ localhost:}\allowbreak{}\texttt{8090/}\allowbreak{}\texttt{health} now fails with a connection error: nothing is listening any more.

\codeneed{16}

Ask Terraform what it thinks:

\codeneed{13}\begin{codeblock}

\begin{Shaded}
\begin{Highlighting}[]
\ExtensionTok{terraform}\NormalTok{ plan}
\end{Highlighting}
\end{Shaded}

\end{codeblock}

\codeneed{11}\begin{codeblock}
\begin{Verbatim}[formatcom=\outputstyle]
Note: Objects have changed outside of Terraform
…
  # docker_container.db has been deleted
…
  # docker_network.lab has been deleted
…
  # docker_container.app will be created
  # docker_container.db will be created
  # docker_network.lab will be created
…
Plan: 3 to add, 0 to change, 0 to destroy.
\end{Verbatim}
\end{codeblock}

\codeneed{7}

Terraform compared its state with reality, found what was missing, and plans to put back exactly that. Apply it:

\codeneed{4}\begin{codeblock}

\begin{Shaded}
\begin{Highlighting}[]
\BuiltInTok{time}\NormalTok{ terraform apply}
\end{Highlighting}
\end{Shaded}

\end{codeblock}

\codeneed{2}\begin{codeblock}
\begin{Verbatim}[formatcom=\outputstyle]
…
Apply complete! Resources: 3 added, 0 changed, 0 destroyed.
\end{Verbatim}
\end{codeblock}

\codeneed{6}

About 14 seconds later:

\codeneed{3}\begin{codeblock}

\begin{Shaded}
\begin{Highlighting}[]
\ExtensionTok{curl}\NormalTok{ localhost:8090/visits}
\end{Highlighting}
\end{Shaded}

\end{codeblock}

\codeneed{1}\begin{codeblock}
\begin{Verbatim}[formatcom=\outputstyle]
{"visits":4}
\end{Verbatim}
\end{codeblock}

Everything is back, and so is the data, because the volume was never deleted. That is the answer to the chapter's question: yes, in seconds, provided the data survives.

\codeneed{13}

\section{Step 5: Destroy and Recreate from Nothing}\label{11-lab-rebuild-everything-from-code:step-5-destroy-and-recreate-from-nothing}

Now remove everything Terraform manages, volume included, answering \texttt{yes}:

\codeneed{7}\begin{codeblock}

\begin{Shaded}
\begin{Highlighting}[]
\ExtensionTok{terraform}\NormalTok{ destroy}
\end{Highlighting}
\end{Shaded}

\end{codeblock}

\codeneed{5}\begin{codeblock}
\begin{Verbatim}[formatcom=\outputstyle]
Plan: 0 to add, 0 to change, 6 to destroy.
…
docker_volume.db_data: Destruction complete after 2s
docker_network.lab: Destruction complete after 2s
Destroy complete! Resources: 6 destroyed.
\end{Verbatim}
\end{codeblock}

\codeneed{10}

Build it again from nothing:

\codeneed{7}\begin{codeblock}

\begin{Shaded}
\begin{Highlighting}[]
\ExtensionTok{terraform}\NormalTok{ apply}
\ExtensionTok{curl}\NormalTok{ localhost:8090/visits}
\end{Highlighting}
\end{Shaded}

\end{codeblock}

\codeneed{4}\begin{codeblock}
\begin{Verbatim}[formatcom=\outputstyle]
…
Apply complete! Resources: 6 added, 0 changed, 0 destroyed.
…
{"visits":1}
\end{Verbatim}
\end{codeblock}

The whole system is back in about 12 seconds, but the count starts again at 1. The code described an empty database, and that is what you got.

\begin{callout}{warning}

Infrastructure as code rebuilds infrastructure, not data. Data needs backups, stored somewhere else, and restores that you have actually tested. In a real cloud, managed databases take the backups for you; your job is to switch them on and try a restore.

\end{callout}

\codeneed{13}

\section{Step 6: Protect the Data from Yourself}\label{11-lab-rebuild-everything-from-code:step-6-protect-the-data-from-yourself}

Terraform can refuse to destroy a resource. Add a \texttt{lifecycle} block to the volume:

\codeneed{7}\begin{codeblock}

\begin{Shaded}
\begin{Highlighting}[]
\NormalTok{resource "docker\_volume" "db\_data" \{}
\NormalTok{  name = "tinyapp{-}db{-}data"}

\NormalTok{  lifecycle \{}
\NormalTok{    prevent\_destroy = true}
\NormalTok{  \}}
\NormalTok{\}}
\end{Highlighting}
\end{Shaded}

\end{codeblock}

\codeneed{8}

and try \texttt{terraform\ destroy} again:

\codeneed{5}\begin{codeblock}
\begin{Verbatim}[formatcom=\outputstyle]
Error: Instance cannot be destroyed
…
Resource docker_volume.db_data has lifecycle.prevent_destroy set, but the
plan calls for this resource to be destroyed.
…
\end{Verbatim}
\end{codeblock}

Nothing was destroyed. Use this for anything whose loss would hurt: databases, buckets, volumes.

\codeneed{9}

\section{Step 7: See Where the Password Went}\label{11-lab-rebuild-everything-from-code:step-7-see-where-the-password-went}

The password never appeared in \texttt{main.tf} or on screen. But look in the state file:

\codeneed{3}\begin{codeblock}

\begin{Shaded}
\begin{Highlighting}[]
\FunctionTok{grep} \AttributeTok{{-}c}\NormalTok{ change{-}me{-}locally terraform.tfstate}
\end{Highlighting}
\end{Shaded}

\end{codeblock}

\codeneed{1}\begin{codeblock}
\begin{Verbatim}[formatcom=\outputstyle]
2
\end{Verbatim}
\end{codeblock}

It is stored there twice, in plain text. This is why state must never go into git, and why teams keep it in encrypted, access-controlled remote state.

To finish, remove the \texttt{lifecycle} block, then run \texttt{terraform\ destroy} and answer \texttt{yes}.

\section{Summary, Key Terms, and Review Questions}\label{11-lab-rebuild-everything-from-code:summary-key-terms-and-review-questions}

\subsection*{Summary}\label{11-lab-rebuild-everything-from-code:summary}

\begin{itemize}
\tightlist
\item
  One Terraform configuration can describe a whole system: images, containers, a network, and a volume, with references deciding the order.
\item
  Variables take inputs such as passwords from \texttt{TF\_VAR\_} environment variables; outputs print useful values.
\item
  A change replaces only what it must. Data kept in a volume or database survives app replacements.
\item
  When things vanish, \texttt{plan} detects it and \texttt{apply} rebuilds exactly what is missing, in seconds.
\item
  Code rebuilds infrastructure, not data. Back data up, and use \texttt{prevent\_destroy} for anything precious.
\item
  State stores secrets in plain text; keep it private.
\end{itemize}

\Needspace{9\baselineskip}

\subsection*{Key Terms}\label{11-lab-rebuild-everything-from-code:key-terms}

\begin{booktable}
\begin{longtable}{>{\raggedright\arraybackslash}p{\dimexpr 0.3582\linewidth-1.2836\tabcolsep\relax}>{\raggedright\arraybackslash}p{\dimexpr 0.6418\linewidth-0.7164\tabcolsep\relax}}
\tabletop
\rowcolor{tablehead}\thead{Term} & \thead{Meaning} \\
\tablemid
\endfirsthead
\tabletop
\rowcolor{tablehead}\thead{Term} & \thead{Meaning} \\
\tablemid
\endhead
\tablebottom
\endlastfoot
Variable (Terraform) & An input to a Terraform configuration \\
Output (Terraform) & A value Terraform prints after applying \\
\end{longtable}
\end{booktable}

\subsection*{Review Questions}\label{11-lab-rebuild-everything-from-code:review-questions}

\begin{enumerate}
\def\labelenumi{\arabic{enumi}.}
\tightlist
\item
  How does Terraform know to create the database before the app?
\item
  After the ``server died'', how did Terraform decide what to create, and why was the data still there?
\item
  What does \texttt{prevent\_destroy} protect against, and what doesn't it protect against?
\item
  Where did the database password end up, even though it was never in a file?
\end{enumerate}

\begin{understandbox}

\textbf{You understand this when} you can delete every container and network of a system without worrying, because you know exactly what one \texttt{terraform\ apply} will bring back, and what it won't.

\end{understandbox}
